\BeleskeID{02-s048b}
% element: 02-s048b:d01

Zamenu troulaznog NI kola pratimo preko međusignala:
\[
u=\overline{XY},\qquad v=\overline{uu}=XY,\qquad
w=\overline{vZ}=\overline{XYZ}.
\]
NI kolo sa spojenim ulazima služi kao invertor. Ako koristimo samo dvoulazna NI kola, ovako dobijamo tri kola za zamenu troulaznog NI. Isti princip koristimo i za četiri ulazne inverzije. Preostala dva NI kola formiraju negirani prvi proizvod i završni izlaz. Sa devet kola u tri četvorostruka kućišta tri kola ostaju neupotrebljena.

Ako međuinverziju \(v=\bar u\) preuzme raspoloživi invertor iz posebnog invertorskog čipa, sama funkcija koristi četiri NI kola: jedno za prvi proizvod, dva u grani koja zamenjuje troulazno NI i jedno završno. Posle četiri ulazne inverzije i ove međuinverzije u šestostrukom invertorskom kućištu ostaje još jedan slobodan invertor. Dodatni nivo na putanji \(X,Y\) treba uzeti u obzir pri proceni kašnjenja.

Pri nabavci upoređujemo stvarnu cenu dva različita kućišta sa cenom tri kućišta iste vrste, uključujući mogući količinski popust. Broj čipova sam po sebi ne određuje ukupnu cenu.
